\section{Origen in Tradition: Memorable Claims under Review}
\label{sec:origen-tradition}

Origen did not acquire a stable saint's Life with an authorized cult. His
memorable traditions grew instead from Eusebius's admiring biography, hostile
doctrinal dossiers, shorthand histories of theology, and modern recovery. The
most vivid claims need neither ridicule nor automatic repetition.

\begingroup
\footnotesize
\setlength{\tabcolsep}{3.5pt}
\begin{longtable}{>{\raggedright\arraybackslash}p{0.20\linewidth}
                  >{\raggedright\arraybackslash}p{0.26\linewidth}
                  >{\raggedright\arraybackslash}p{0.23\linewidth}
                  >{\raggedright\arraybackslash}p{0.23\linewidth}}
\toprule
\textbf{Memorable claim} & \textbf{Earliest controlled witness} &
\textbf{Historical status} & \textbf{Responsible use}\\
\midrule
\endhead
Origen was a pagan Greek who converted as an adult. & Porphyry, quoted by
Eusebius at \emph{Ecclesiastical History} 6.19, says Origen moved from Greek
life to Christian teaching; Eusebius replies that Origen was Christian from
childhood and has already narrated Leonides. & The witnesses conflict. The
martyr father and scriptural childhood are integral to the earlier sustained
biographical account; Porphyry is hostile but valuable and may have understood
Origen through a conversion stereotype or confused biographical detail. &
Preserve the conflict explicitly. ``Greek'' describes education and culture in
the passage, not a modern ethnicity proved by a civil record.\\
Adamantius means that Origen was literally ``made of steel.'' & Eusebius knows
the additional name; later writers explain it through indefatigable labor. &
The name is early; whether birth name, sobriquet, and exact original force are
unresolved. & It can symbolize endurance only when interpretation is not
presented as a registered surname or physical fact.\\
Origen founded a university called the Catechetical School of Alexandria. &
Eusebius 6.3--6, 6.15 describes teaching, Demetrius's commission, and a later
division of pupils; no charter survives. & A real episcopally connected and
evolving teaching activity; modern institutional form unproved. & ``School''
may name teachers, pupils, curriculum, and succession without importing
departments, degrees, buildings, or uninterrupted corporate identity.\\
Clement was certainly Origen's long-term personal master. & Eusebius 6.6 calls
Origen a pupil of Clement; the historical overlap is plausible. & Early
received relation, but duration, setting, and curriculum lack independent
control. & Speak of probable direct influence alongside the much clearer
literary and Alexandrian inheritance.\\
Origen studied with the same Ammonius Saccas who taught Plotinus. & Porphyry's
extract names Ammonius; identification with Plotinus's teacher is later
reconstruction. & Possible K, not established; ancient people shared the name,
and Eusebius also disputes Porphyry's description of Ammonius. & Use only to
establish serious philosophical formation, not a documented classroom lineage
from Ammonius through Origen to Christian Platonism.\\
Origen's mother hid his clothes to stop his martyrdom. & Eusebius 6.2, written
about a century after the reported event. & Specific E biographical tradition,
not independently corroborated. & It truthfully depicts remembered maternal
agency and adolescent zeal if not embellished into eyewitness dialogue.\\
Origen castrated himself because of Matthew 19:12. & Eusebius 6.8 gives the
specific act, motives, and Demetrius's later use; Origen's surviving Matthew
commentary rejects literal enactment. & Early, disputed report. No admission,
independent contemporary witness, or bodily evidence survives. & Keep as a
serious unresolved claim and warning about ascetic literalism, not a comic
anecdote, medical certainty, or fact to conceal.\\
Demetrius expelled Origen solely from jealousy. & Eusebius gives jealousy as
the moral explanation; original Alexandrian acts are lost. & Rupture secure;
single motive unproved. & Include jurisdiction, irregular ordination, doctrine,
ascetic status, and personality as possible factors without inventing their
weight.\\
Two Alexandrian synods precisely tried and deposed him for enumerated heresies.
& Photius, ninth century, \emph{Bibliotheca} codex 118, summarizes Pamphilus's
lost apology and distinguishes two assemblies. & May preserve early material,
but it is late and lacks the
acts, charges, votes, and Origen's complete defense. & State two synods as a
late dossier reconstruction, not a contemporary transcript; do not collapse
the rupture into a modern heresy trial.\\
Ambrose was definitively a Valentinian---or definitively a Marcionite---before
Origen converted him. & Eusebius 6.18/6.23 gives the Valentinian affiliation;
Jerome, \emph{On Illustrious Men} 56, calls him Marcionite. & Conversion and
patronage converge; prior affiliation conflicts. & Preserve both labels rather
than synthesizing a sectarian biography not supplied by either witness.\\
Julia Mamaea became Christian or was baptized by Origen. & Eusebius 6.21 says
she summoned Origen and heard Christian teaching, while praising her piety. &
Audience and instruction E; baptism and formal affiliation unproved. & The
episode shows Christian teaching reaching imperial proximity, not a secret
Christian court.\\
Philip the Arab was baptized by Origen, the first Christian emperor. & Eusebius
6.34 preserves a separate Easter story about an unnamed emperor, and 6.36 notes
Origen's letters to Philip and Severa; later traditions combine claims. &
Letters likely; baptism by Origen and identification of the penitent emperor
are unproved in the controlled record. & Do not make correspondence a
sacramental dossier or settle Philip's religion by harmonization.\\
The Hexapla was exactly six columns throughout and the first modern critical
Bible. & Eusebius 6.16--17 and surviving fragments establish a comparative
project with varying extra versions. & Monumental textual scholarship secure;
physical form, stages, and goals debated. & Explain the ancient philological
project on its terms and identify ``six-column'' or ``critical edition'' as
shorthand.\\
Origen wrote exactly six thousand books. & Ancient and later work lists count
titles, books within works, homilies, letters, and scholia differently. & Huge
output secure; exact popular totals legendary or methodologically unstable. &
Name the surviving work and transmission instead of using a prodigious number
as biographical proof.\\
Origen taught that all souls, necessarily including the devil, will certainly
be saved. & Restoration language occurs in authentic works; other passages
stress freedom and judgment. Rufinus, fragments, and later Origenist dossiers
differ. & Universal restoration is a real Origenian trajectory; its modal
certainty, mechanics, recurrence, and inclusion of the devil remain disputed.
& Neither erase the trajectory nor quote the maximal later system as a settled
sentence by Origen. Catholic teaching supplies a reception boundary.\\
Origen denied bodily resurrection. & His works confess resurrection and argue
from 1 Corinthians 15 for transformed identity; ancient critics challenge his
account of bodily continuity. & Simple denial is false; some formulations are
genuinely problematic and not identical with later received teaching. & Keep
the confession, the metaphysical account, the criticism, and later doctrinal
boundary distinct.\\
Origen's allegory means that biblical events never happened. & Commentaries
combine variants, grammar, history, moral address, and spiritual senses, while
\emph{On First Principles} sometimes identifies narrative impossibilities. &
False as a general rule; the balance varies by text and passage. & Test each
reading rather than using ``allegorist'' as praise or dismissal.\\
Origen was condemned as a heretic during his lifetime everywhere in the Church.
& Alexandria repudiated his ministry; Palestinian and other bishops continued
to receive and consult him. Exact Alexandrian acts are lost. & False
universalization of a real local rupture. & Name jurisdiction, action, and
divided reception.\\
Constantinople II either condemned every sentence Origen wrote or never
condemned Origen at all. & The council's received anathema 11 names Origen and
impious writings; the famous fifteen anti-Origenist anathemas are not in the
formal acts and are excluded by Tanner from the conciliar decrees. & Both
absolutes are false. Relation of the fifteen to pre-conciliar deliberation
remains debated. & State the received act's actual wording and distinguish
author, authentic works, later systems, the 543 action, and disputed textual
attachments.\\
Origen was a canonized saint and martyr whose feast was later suppressed. & No
controlled canonization, universal ancient cult, stable feast, or execution
record was located; Eusebius says he survived torture. & Confessor secure;
saint, martyr by death, and suppressed-feast narrative unproved. & Admiration
and private honor should not be represented as a competent liturgical act.\\
Benedict XVI restored Origen to official Catholic authority or canonized him in
2007. & Two
official general audiences praise Origen as exegete and teacher. & Important M
catechetical reception; no canonization or new conciliar act. & Read
the praise in its actual genre and scope.\\
Any elegant spiritual sentence circulating under Origen's name is authentic. &
The corpus survives through translations, fragments, catenae, and later
anthologies; online lists commonly omit work and locus. & Unsupported until
work, section, language, edition, and translator are identified. & Exact
attribution matters especially for an author whose translators admitted
correction.\\
\bottomrule
\end{longtable}
\endgroup

\subsection{Why doubtful stories still matter}

The clothes, the bodily act, the iron endurance, and the condemned name have
become miniature explanations of Origen. Each makes a genuine feature visible:
maternal rescue, ascetic danger, intellectual labor, fidelity under pain, or
speculation's risk. Each becomes false when it replaces the documentary scale
of a life.

Tradition history is therefore not a disposal appendix. It asks why Christians
needed Origen to be the literalist who harmed his body, the allegorist who
escaped bodies, the universalist who abolished judgment, or the master secretly
restored by a pope. The contradictions reveal continuing arguments over how
Scripture, reason, authority, hope, and error belong together.
